% Complete Section 8.2, original PDF 108-109 / printed 98-99.
% block: variable-coefficients-heading page:108 type:section
\subsection{{Variable coefficients}}\label{sec:variable-coefficients}

% block: weighted-form-inner-product pages:108,109 type:ordinary-paragraph
{Next we begin to move beyond the Hodge Laplacian. A simple variation is to change the inner product on one or more of the base Hilbert spaces $W^k=L^2\Lambda^k$ from the standard $L^2$ inner product to a weighted inner product:}
\[
\langle\omega,\mu\rangle=\int_\Omega\langle A\omega,\mu\rangle\operatorname{vol},\qquad\omega,\mu\in L^2\Lambda^k(\Omega).
\]
{Here $A$ is a function on $\Omega$ for which $A(x)$ is a symmetric positive definite operator on $\operatorname{Alt}^kT_x\Omega$ for each $x\in\Omega$ (e.g., if $\Omega$ is a domain in $\mathbb R^n$, $A$ can be viewed as a function from $\Omega$ to the space of symmetric positive definite matrices of size $\binom nk$). Assuming an upper bound on $A$ and a positive lower bound on its least eigenvalue, the weighted inner product on $W^k$ is equivalent to the usual one, and so the new Hilbert complex inherits all essential properties as the old, such as having closed ranges. The abstract Hodge Laplacian of the complex, however, is different. Leaving the general case to the reader, consider the simple example of the final segment of the de Rham complex on a domain in $\mathbb R^n$. We take as the base spaces the weighted space $L^2\Lambda^{n-1}(\Omega,A)$ described above, where $A$ is an $n\times n$ symmetric positive definite matrix-valued function, and leave the rest of the complex untouched, so the relevant segment is}
\begin{equation}\label{eq:weighted-divergence-complex}
L^2\Lambda^{n-1}(\Omega,A)\xrightarrow{(\operatorname{div},H(\operatorname{div}))}L^2\Lambda^n(\Omega)\to0.
\end{equation}

% block: variable-coefficient-mixed-problem page:109 type:ordinary-paragraph
{The weak form of the abstract Hodge Laplace problem for this complex then seeks $\sigma\in H(\operatorname{div})$, $u\in L^2$ such that}
\[
\begin{aligned}
\langle A\sigma,\tau\rangle-\langle u,\operatorname{div}\tau\rangle&=0,&\tau&\in H(\operatorname{div}),\\
\langle\operatorname{div}\sigma,v\rangle&=\langle f,v\rangle,&v&\in L^2.
\end{aligned}
\]
{This is the standard weak formulation corresponding to the problem}
\[
A\sigma+\operatorname{grad}u=0,\qquad\operatorname{div}\sigma=f,
\]
{or, equivalently, $-\operatorname{div}C\operatorname{grad}u=f$, $\sigma=-C\operatorname{grad}u$, where $C=A^{-1}$, together with the Dirichlet boundary condition $u=0$ on $\partial\Omega$ which, in this weak formulation, is a natural boundary condition. In short, this variable coefficient partial differential operator is again the abstract Hodge Laplacian of a Hilbert complex, albeit a slightly different Hilbert complex than the $L^2$ de Rham complex. It may be discretized with the same finite element spaces as in the case where $A$ is the identity, and the analysis carries over without additional considerations. Many other elliptic partial differential equations with variable coefficients can be treated in a similar fashion.}
